\documentclass{../note}

\begin{document}

\begin{zadanie}{1}
Dwa prostokąty na płaszczyźnie są w relacji $r$ wtedy i tylko wtedy, gdy istnieje przesunięcie przekształcające jeden na drugi.

Ile klas abstrakcji ma relacja $r$ i jakie są ich moce?
\end{zadanie}
\begin{rozwiazanie}
Każdy prostokąt jest wyznaczony swoim lewym dolnym wierzchołkiem, czyli parą współrzędnych $(x, y) \in \R^2$, oraz szerokością i wysokością $w, h \in \R_+$. Wtedy dwa prostokąty są w relacji, czyli jeden jest przesunięciem drugiego, gdy mają równą wysokość i szerokość, czyli wysokość i szerokość wyznaczają klasę abstrakcji. Zatem jest tyle klas abstrakcji, ile par $(w, h)$, czyli $|\R_+^2|$. $\R_+$ zawiera odcinek $(0, 1)$, o którym wiemy, że ma moc $\mathfrak{c}$, oraz jest podzbiorem $\R$, które ma moc $\mathfrak{c}$, więc $|\R_+| = |\R|$. Zatem zbiór klas abstrakcji ma moc $|\R_+^2| = |\R^2| = \mathfrak{c}$.\\
Wszystkie prostokąty z jednej klasy abstrakcji są wyznaczone przez pozycję, czyli przez $(x, y)$. Zatem moc każdej klasy abstrakcji to $|\R^2| = \mathfrak{c}$.
\end{rozwiazanie}

\end{document}